\section{Introduction}

Streamflow prediction plays a critical role in water resource management, flood control, and disaster mitigation, as well as basin operation and regulation\cite{GAO2026103148}. Under intensifying hydro-climatic uncertainty\cite{immerzeel2010climate}, such predictions increasingly support high-stakes decision-making processes, including water allocation, flood-risk governance, and adaptive basin management\cite{milly2020colorado}. In these contexts, model outputs not only guide engineering operations but may also influence regional or even transboundary water resource negotiations. As a result, model trustworthiness has become as important as predictive accuracy.

In recent years, machine learning and deep learning models have achieved significant improvements in streamflow prediction, demonstrating strong capability in capturing complex nonlinear relationships and temporal dependencies\cite{apak2026multi,ZIPPER2026134909}. However, these models are often criticized as ``black boxes''\cite{Núñez2023_XAI}, whose internal decision-making processes are difficult to interpret intuitively, thereby limiting the transparency and reliability of predictions\cite{WangShap2024}. In non-stationary or extrapolation scenarios, even highly accurate models may produce unreliable predictions if their internal logic is not understood, potentially misleading basin management decisions. Therefore, improving interpretability while maintaining predictive performance has become a central challenge in hydrological modeling\cite{MardShap2022}.

To mitigate the black-box nature of data-driven models, researchers have increasingly incorporated hydrological process knowledge into neural networks to enhance physical consistency\cite{SCHMID2026124819}. A representative approach is Physics-Informed Neural Networks (PINNs), which embed governing physical equations into the loss function as constraints, enabling models to satisfy both observational data and physical laws\cite{SHI2025_PINN,LIU2026_PINN}. For example, Fang et al.\cite{FANG2025_PINN} proposed the SWAN model, a PINN-based framework that integrates saturated streamflow process knowledge into the neural network architecture. While such approaches improve physical plausibility to some extent, physical consistency does not inherently guarantee interpretability\cite{Roscher2020_Explainable}. The internal mechanisms by which models utilize input variables\cite{NASER2026_Inrter} and represent hydrological processes remain largely unclear\cite{Krishnapriyan2021_PINN}.

Current mainstream interpretability methods primarily rely on post-hoc analysis techniques, such as LIME\cite{Ribeiro2016_LIME_Origin} and SHAP\cite{lundberg2017_Shap_Origin}. These approaches interpret complex models by constructing local surrogate approximations or decomposing predictions based on cooperative game theory\cite{Clark2025_Shap}, thereby quantifying the contribution of input features to model outputs. In this way, they enhance the transparency of otherwise black-box models to a certain extent\cite{Wagan2026_LIME}. In streamflow prediction studies, such methods provide intuitive tools for analyzing the influence of key variables, including precipitation, evapotranspiration, and soil moisture, and have consequently become among the most widely adopted interpretability techniques.

However, the explanations produced by these methods generally lack explicit verifiability\cite{NEURIPS2021_shap_residuals}. Unlike model predictions, which can be directly evaluated against observed streamflow, feature attributions do not have an objective ground truth and are therefore difficult to validate through independent experiments\cite{TAKEFUJI2026_Shap_problem}. Existing studies have shown that even when a model achieves high predictive accuracy, its explanations may partly reflect complex statistical fitting or spurious correlations\cite{e24050687_Shap_problem}. As a result, relying solely on feature attribution makes it difficult to determine whether a model has genuinely learned physically meaningful relationships\cite{Benchmarkingofa_PhysicallyBasedHydrologicModel}, which limits its credibility in high-stakes decision-making contexts.

Furthermore, most existing studies remain at the level of variable-level analysis based on SHAP or LIME\cite{ZHOU2026103254,XIA2026103232}, lacking a systematic evaluation of model behavior. While these methods provide useful insights into the relative importance of input variables, they do not assess how models respond under changing conditions\cite{10.2166/wcc.2017.149}. In hydrological systems, however, understanding model behavior is essential. A physically reliable model is expected to produce directionally consistent and hydrologically reasonable responses\cite{hess-30-1077-2026} to variations in key drivers, such as precipitation.

Specifically, limited attention has been given to how models behave under input perturbations or counterfactual scenarios, resulting in a lack of systematic sensitivity analysis\cite{simic2025perturbation}. For example, when precipitation magnitude increases, streamflow volume should correspondingly increase in a stable and consistent manner\cite{nearing2021role}. If a model fails to exhibit such reasonable response patterns under varying input magnitudes, its explanations—even if statistically consistent—cannot be regarded as process-relevant. Therefore, interpretability in hydrological modeling should move \textbf{\emph{beyond feature attribution}} and incorporate behavioral validation, such as sensitivity analysis and counterfactual experiments\cite{ZHANG2025123192}, to establish consistency between explanation results, model behavior, and underlying hydrological processes\cite{read2019process}.

In addition, existing interpretability approaches are largely confined to individual models and lack systematic evaluation\cite{pmlr-v235-wang24j} of whether the learned knowledge is transferable. If the key feature relationships or internal structures identified in one model cannot be reproduced or improve performance in other model frameworks, such explanations are more likely to reflect model-specific parameterization rather than generalizable hydrological knowledge. From a scientific perspective, meaningful explanations should hold across different model architectures and be validated beyond a single model\cite{li2022transferability}. Therefore, cross-model validation is a crucial step in determining whether the extracted knowledge possesses true scientific value.

In summary, current interpretability studies in deep learning-based streamflow prediction face three major challenges: (1) \textbf{lack of verifiability}, making it difficult to quantitatively validate explanation results; (2) \textbf{lack of behavior-level consistency evaluation}, particularly the absence of sensitivity analysis to assess whether model responses align with hydrological principles; and (3) \textbf{lack of transferability assessment}, leaving it unclear whether the learned knowledge is model-independent and generalizable. These limitations collectively constrain the scientific interpretability and practical reliability of existing approaches.Notably, other fields have increasingly shifted toward rigorous and verification-oriented interpretability paradigms, such systematic efforts remain limited in hydrological modeling.

To address the three major challenges, this study proposes a novel framework, Sparse Transformer-based Explainable Hydrology (STEH), which enables full-chain interpretability from inputs through internal model structure to outputs. Specifically, (1) A sparse circuit discovery mechanism is introduced via learnable gating across input features, attention heads, and neurons, allowing the model to identify compact and functionally meaningful computational pathways (e.g., evaporation, precipitation–runoff transformation); (2) circuit-level ablation experiments are conducted to quantitatively verify the contribution of individual components; (3) internal dependency structures are analyzed to reveal information flow and align model representations with hydrological processes; (4) sensitivity and counterfactual analyses are performed to evaluate whether model responses are physically consistent; and (5) cross-model transfer experiments are designed by incorporating the extracted structural knowledge into a Random Forest benchmark, thereby assessing the transferability of learned knowledge.

The main contribution of this study lies not only in the development of an interpretable deep learning framework, but also in establishing a systematic validation pathway for interpretability in hydrological modeling. By integrating structural analysis, behavioral verification, and cross-model evaluation, this work advances the understanding of interpretability from merely explaining model outputs to rigorously assessing whether such explanations are reliable, physically meaningful, and transferable. In doing so, it provides preliminary insights into a fundamental question: \textbf{\emph{what kind of hydrological knowledge is learned, and whether it is verifiable and transferable?}}